\section{Background and Related Work}
\label{sec:background}

\subsection{PhaseNet and Deep Learning for Seismology}

% One paragraph. What PhaseNet is (Zhu & Beroza 2019): U-Net CNN, three-component waveforms → P/S probability traces + confidence.
% Mention SeisBench (Woollam et al. 2022) for evaluation interface.
% Note the trend: models trained on well-instrumented networks deployed globally → distribution shift.
% Use Sera's description of the architecture but keep it concise.

PhaseNet \citep{zhu2019phasenet} is a U-Net convolutional neural network that takes three-component seismic waveforms as input and outputs continuous probability traces for P-wave, S-wave, and noise classes. Picks are extracted via peak detection with configurable confidence thresholds. The model was trained on approximately 780{,}000 waveforms from the Northern California Seismic Network and has become one of the most widely-used tools for automated phase picking, accessible through the standardised SeisBench interface \citep{woollam2022seisbench}. It is increasingly deployed to data-sparse regions globally---the distribution shift scenario we study here.

\subsection{Distribution Shift and the Limits of Model Self-Assessment}
\label{sec:bg_shift}

% Two paragraphs.
% Para 1: Neural network miscalibration (Guo et al. 2017), degradation under shift (Ovadia et al. 2019).
% Core module concept: epistemic uncertainty should increase OOD → natural deferral signal.
% We test this for PhaseNet and find it fails.
% Para 2: Selective prediction history (Chow 1970, Geifman & El-Yaniv 2017), learning to defer (Madras et al. 2018, Mozannar & Sontag 2020).
% Key distinction: most assume (a) informative uncertainty or (b) labelled deployment data. Neither holds here.
% Our gap: external environmental context instead of model self-assessment.

Neural network confidence scores are often poorly calibrated even on in-distribution data \citep{guo2017calibration}, and calibration degrades further under distribution shift \citep{ovadia2019uncertainty}. A core principle of uncertainty quantification is that epistemic uncertainty should increase in unfamiliar regions of input space, providing a natural signal for selective prediction or deferral. We test whether this holds for PhaseNet and find that it does not: confidence scores carry essentially no information about whether the model is operating in- or out-of-distribution, with P-wave AUROC virtually identical between California (0.619) and Kazakhstan (0.616). This motivates the search for an external reliability signal.

Selective prediction and the reject option have a long history \citep{chow1970optimum,geifman2017selective}. The learning-to-defer framework \citep{madras2018predict,mozannar2020consistent} formalises the decision of when to defer to a human expert, optimising a joint objective over automated accuracy and expert review cost. A key distinction for our setting is that most deferral frameworks assume either that model uncertainty is informative or that a learned rejector can be trained on labelled deployment-domain data. Neither holds here: confidence is uninformative, and labelled Kazakhstan data is scarce. Our framework addresses this gap by using \emph{external environmental context} rather than model self-assessment.

\subsection{HMMs and Seismic Regime Modelling}
\label{sec:bg_hmm}

% One paragraph. Omori's law (1894), ETAS (Ogata 1988), the Q→A→D cycle.
% HMM is domain-informed, not arbitrary. Connects to module emphasis on incorporating domain knowledge.

The seismic cycle of quiet background $\to$ mainshock $\to$ aftershock decay is well-established in seismology \citep{omori1894aftershocks,ogata1988etas}. Aftershock rates follow Omori's law, decaying as a power law with time since the mainshock \citep{utsu1995centenary}. This Quiet $\to$ Active $\to$ Decaying pattern maps directly onto Hidden Markov Model assumptions: hidden states generating observable count data through state-dependent emission distributions \citep{rabiner1989hmm}. Our use of an HMM is therefore not an arbitrary modelling choice but a domain-informed one---we encode established seismological structure into the monitoring framework rather than learning regime dynamics from scratch.
